\chapter{Audio Effects and Synthesis}

So far we have spent our energy \emph{analyzing} sound: pulling it apart into
frequencies, smoothing it, measuring its features. This chapter flips the arrow.
We are going to \emph{make} and \emph{shape} sound. An echo, a chorus, a fuzzy
guitar tone, the breathy attack of a flute, the metallic bite of an FM bell ---
every one of these is a signal-processing recipe, and once you can read the
recipe you can hear the math.

The wonderful thing is that you already own all the tools. An audio effect is
almost always a filter, a delay, or a nonlinearity --- and we have studied
filters and delays at length. A synthesizer is almost always a small bank of
oscillators run through a filter and an amplitude envelope --- and an oscillator
is just the rotating phasor of the early chapters sampled in time. So this
chapter is less about new theory and more about \emph{seeing} how a handful of
familiar pieces, wired together, produce the entire palette of electronic and
processed sound.

Our plan: first effects, organized by the operation underneath them ---
amplitude (gain, distortion), spectrum (EQ), and delay (echo, reverb, flanger,
chorus). Then synthesis, organized by how the timbre is built --- adding
sinusoids up (additive), carving them away from a rich source (subtractive),
modulating one oscillator's phase with another (FM), reading a stored cycle
(wavetable), or simulating a vibrating string with a feedback delay loop
(physical modeling). Throughout we keep the geometry in view: a timbre is a shape in the
frequency domain, an envelope is a shape in the time domain, and an effect is a
transformation of one or both. Short code sketches make every idea concrete, and
the chapter closes by connecting these synthesis engines to the wrist controller.

\chapterline

\section{The three knobs: amplitude, spectrum, time}
\label{ch28:three-knobs}

Strip away the marketing names on a guitar pedal or a plugin and you find only
three things being manipulated.

\begin{itemize}
  \item \textbf{Amplitude} --- how loud each sample is. Scaling it by a constant
        is gain. Scaling it by a slowly varying curve is tremolo or an envelope.
        Passing it through a \emph{nonlinear} curve is distortion.
  \item \textbf{Spectrum} --- how energy is distributed across frequency.
        Reshaping it is equalization (EQ) and, more generally, filtering.
  \item \textbf{Time} --- when each sample is heard. Copying the signal to a
        later time and adding it back is delay, and delay is the seed of echo,
        reverb, flanging, and chorus.
\end{itemize}

\begin{intuitionbox}
Almost every effect is a combination of these three, plus one twist:
\emph{modulation}, where a slow control signal (a low-frequency oscillator, an
envelope, or a sensor reading from your wrist) drives one of the three knobs over
time. A static filter is an EQ; a filter whose cutoff sweeps is a wah pedal. A
fixed delay is an echo; a delay whose length wobbles is a flanger. The
\emph{movement} is what makes an effect feel alive, and that movement is exactly
where a gesture controller will plug in.
\end{intuitionbox}

Two of these knobs we can already turn cold. Gain is multiplication by a scalar
$g$: $y[n]=g\,x[n]$. EQ is the LTI filtering of the earlier chapters: choose a
frequency response $H(e^{\jj\omega})$ and convolve. The two genuinely new ideas
in this chapter are \emph{nonlinearity} (distortion, and the heart of FM
synthesis) and \emph{structured delay} (the comb and all-pass building blocks of
spatial effects). We take them in turn.

\section{Amplitude effects: gain, tremolo, and distortion}
\label{ch28:amplitude}

\subsection{Gain and tremolo}

The simplest effect imaginable is a volume knob:
\begin{equation}
\label{ch28:gain}
y[n] \;=\; g\,x[n].
\end{equation}
This is linear and time-invariant, so it changes nothing about the spectrum's
\emph{shape} --- every frequency is scaled by the same $g$. Make $g$ vary slowly
in time and you get \textbf{tremolo}:
\begin{equation}
\label{ch28:tremolo}
y[n] \;=\; \bigl(1 + d\,\cos(2\pi f_m n / f_s)\bigr)\,x[n],
\end{equation}
with modulation depth $d\in[0,1]$ and a modulation rate $f_m$ of a few hertz.
Now $g$ is time-\emph{varying}, so this is no longer LTI; in the spectrum it
splits each tone into sidebands at $\pm f_m$ (amplitude modulation, the same
mechanism we met in Chapter~23). At a few hertz you hear it as a
pulsing in loudness rather than as new pitches
(Figure~\ref{ch28:fig-tremolo}).

\begin{figure}[H]
\centering
\begin{tikzpicture}
\begin{axis}[
  width=12.5cm, height=4.2cm,
  xlabel={$n$ (samples)}, ylabel={amplitude},
  xmin=0, xmax=200, ymin=-1.25, ymax=1.25,
  xtick={0,50,100,150,200}, ytick={-1,0,1},
  axis lines=left, ymajorgrids, grid style={gray!18},
  samples=200, title={\small tremolo: a slow loudness envelope riding on the tone},
]
\addplot[domain=0:200, thick, sigblue] {(1+0.6*cos(deg(2*pi*x/120)))*sin(deg(2*pi*x/12))/1.6};
\addplot[domain=0:200, gray!70, dashed] {(1+0.6*cos(deg(2*pi*x/120)))/1.6};
\addplot[domain=0:200, gray!70, dashed] {-(1+0.6*cos(deg(2*pi*x/120)))/1.6};
\end{axis}
\end{tikzpicture}
\caption{Tremolo: the audio-rate tone (solid) is scaled by a slow loudness curve
(dashed envelope), Eq.~\eqref{ch28:tremolo}. The envelope rises and falls at
$f_m$ a few times per second, so you hear a pulsing volume rather than a change
in pitch.}
\label{ch28:fig-tremolo}
\end{figure}

\subsection{Distortion: bending the line}

Gain multiplies by a number. \textbf{Distortion} sends each sample through a
\emph{curved} function $f(\cdot)$, applied sample by sample:
\begin{equation}
\label{ch28:waveshape}
y[n] \;=\; f\bigl(x[n]\bigr).
\end{equation}
This is a \emph{memoryless nonlinearity} (also called waveshaping): the output at
time $n$ depends only on the input at time $n$, but not linearly. That single
word ``nonlinear'' is doing enormous work. A nonlinear function manufactures
\emph{new frequencies} that were not in the input.

\begin{whybox}
Why does bending the amplitude curve create new frequencies? Feed a pure tone
$x[n]=\cos(\omega n)$ into a polynomial nonlinearity, say $f(x)=x^2$. Then
\[
y[n]=\cos^2(\omega n)=\tfrac12+\tfrac12\cos(2\omega n).
\]
A DC term and a tone at \emph{twice} the frequency appeared --- harmonics that
the input never contained. Cube it and you get a tone at $3\omega$ as well. Any
nonlinear $f$ can be expanded as a power series $f(x)=a_0+a_1x+a_2x^2+\cdots$,
and the $x^k$ term generates energy up to the $k$-th harmonic. Bending the line
sprays the spectrum upward into harmonics. \emph{That harmonic spray is the
sound of distortion}: a clean sine becomes a buzzy, rich tone.
\end{whybox}

The most important shape in practice is a \textbf{saturating} curve --- one that
is nearly straight for small inputs but flattens out for large ones, like an
amplifier running out of headroom. Two common choices:
\begin{align}
\text{soft clip (tanh):}\quad & f(x)=\tanh(k\,x), \label{ch28:tanh}\\
\text{hard clip:}\quad & f(x)=\max\!\bigl(-1,\min(1,k\,x)\bigr). \label{ch28:hardclip}
\end{align}
The gain $k$ is the ``drive.'' Small $k$ stays in the straight middle of the
curve and barely changes the sound; large $k$ pins the waveform against the flat
top and bottom, squaring it off.

\begin{figure}[H]
\centering
\begin{tikzpicture}
\begin{axis}[
  width=7.2cm, height=5.2cm,
  xlabel={input $x$}, ylabel={output $f(x)$},
  xmin=-3, xmax=3, ymin=-1.4, ymax=1.4,
  xtick={-3,-1.5,0,1.5,3}, ytick={-1,0,1},
  axis lines=left, ymajorgrids, grid style={gray!18},
  legend style={at={(0.02,0.98)},anchor=north west,font=\scriptsize,draw=gray!40},
  title={\small saturating curves}, samples=120,
]
\addplot[domain=-3:3, thick, sigblue] {tanh(2*x)};
\addlegendentry{$\tanh(2x)$}
\addplot[domain=-3:3, thick, specred, dashed] {max(-1,min(1,2*x))};
\addlegendentry{hard clip}
\addplot[domain=-3:3, gray, dotted] {x};
\addlegendentry{$y=x$ (clean)}
\end{axis}
\end{tikzpicture}
\hfill
\begin{tikzpicture}
\begin{axis}[
  width=7.2cm, height=5.2cm,
  xlabel={$n$ (samples)}, ylabel={amplitude},
  xmin=0, xmax=60, ymin=-1.4, ymax=1.4,
  xtick={0,20,40,60}, ytick={-1,0,1},
  axis lines=left, ymajorgrids, grid style={gray!18},
  legend style={at={(0.02,0.02)},anchor=south west,font=\scriptsize,draw=gray!40},
  title={\small clean sine vs.\ saturated}, samples=120,
]
\addplot[domain=0:60, gray, dotted, thick] {0.9*sin(deg(2*pi*x/40))};
\addlegendentry{clean}
\addplot[domain=0:60, thick, sigblue] {tanh(4*0.9*sin(deg(2*pi*x/40)))};
\addlegendentry{$\tanh$, drive 4}
\end{axis}
\end{tikzpicture}
\caption{Left: saturating transfer curves. Near the origin they follow $y=x$
(clean); far out they flatten, clipping the peaks. Right: a sine pushed through
$\tanh$ with high drive --- its rounded top and bottom flatten toward a square
wave, which is rich in odd harmonics.}
\label{ch28:fig-distortion}
\end{figure}

\begin{happeningbox}
What is really happening when you ``turn up the gain'' on a distortion pedal is
not that it gets louder --- the saturating curve caps the output amplitude. What
changes is the \emph{spectrum}. As drive increases, the waveform morphs from a
sine toward a square, and the spectrum fills in with harmonics: $3f,5f,7f,\dots$
for a symmetric (odd) curve like $\tanh$, and even harmonics too if the curve is
asymmetric (which is why tube amps, slightly asymmetric, sound ``warmer'' than
symmetric transistor clipping). Shaping amplitude nonlinearly in time is simultaneously
a harmonic-generation process in frequency.
\end{happeningbox}

\begin{trapbox}
\textbf{Distortion aliases.} Because waveshaping creates harmonics far above the
input frequency, those harmonics can easily exceed the Nyquist limit $f_s/2$ and
fold back as inharmonic, dissonant aliases. A clean-sounding software distortion
almost always \emph{oversamples}: upsample by $4\times$--$8\times$, apply the
nonlinearity where there is room for the new harmonics, low-pass filter, then
decimate back. Skip this and a high note run through heavy distortion turns to
gravel.
\end{trapbox}

\begin{lstlisting}[language=Python,caption={Soft-clip distortion with a drive knob.}]
import numpy as np

def distort(x, drive=4.0):
    # tanh soft-clip; normalize so quiet passages are unchanged
    return np.tanh(drive * x) / np.tanh(drive)

# (Production code would oversample around this line to tame aliasing.)
\end{lstlisting}

\section{Spectral effects: EQ as everyday filtering}
\label{ch28:eq}

An \textbf{equalizer} is just a filter (or a bank of filters) you adjust by ear.
There is no new mathematics here --- it is the frequency-response shaping of
Chapters~14 and~16 --- but the vocabulary is worth pinning down because it is how
musicians and engineers talk about spectrum.

\begin{itemize}
  \item \textbf{Shelving filters} lift or cut everything above (high shelf) or
        below (low shelf) a corner frequency by a fixed amount. Think of the
        ``bass'' and ``treble'' knobs on a stereo.
  \item \textbf{Peaking / bell filters} boost or cut a band centered on a chosen
        frequency, with a width set by the quality factor $Q$. A graphic EQ is a
        row of fixed bells; a parametric EQ lets you move each bell's center,
        gain, and $Q$.
  \item \textbf{High-pass / low-pass} remove rumble below or hiss above a cutoff.
\end{itemize}

These are almost universally built as \textbf{biquads} --- second-order IIR
sections, Equation~\eqref{ch28:biquad} --- because two poles and two zeros are
exactly enough to place one resonant bump or notch, and you cascade several to
build a complicated curve.
\begin{equation}
\label{ch28:biquad}
y[n]=b_0x[n]+b_1x[n-1]+b_2x[n-2]-a_1y[n-1]-a_2y[n-2].
\end{equation}
The well-known ``audio EQ cookbook'' formulas give the five coefficients
$\{b_0,b_1,b_2,a_1,a_2\}$ for each EQ type directly from $(f_c, Q, \text{gain})$;
in practice you compute them once whenever a knob moves and run the recurrence per
sample.

\begin{connectionbox}
A biquad's two poles are a complex-conjugate pair sitting at radius $r$ and angle
$\theta=2\pi f_c/f_s$ inside the unit circle, as Chapter~15 showed.
The closer $r$ creeps to $1$, the sharper and taller the resonant peak in the
frequency response --- the same pole-near-the-circle geometry that governs every
resonator, from a plucked string to a vowel formant in the human voice. Studio EQ,
subtractive synth filters, and physical acoustic resonators all share this exact pole-near-the-unit-circle geometry.
\end{connectionbox}

\section{Delay: the seed of echo, reverb, flanging, and chorus}
\label{ch28:delay}

Now the genuinely spatial effects. They all grow from one trivial operation: a
\textbf{delay line}, a buffer that hands back what you put in $D$ samples ago,
i.e.\ $x[n-D]$. The delay length in seconds is $\tau=D/f_s$. What you do with the
delayed copy --- and whether you feed it back --- determines the whole family.

\subsection{Single echo and the feedforward comb filter}

Add one delayed, attenuated copy to the original:
\begin{equation}
\label{ch28:feedforward-comb}
y[n]=x[n]+g\,x[n-D],\qquad 0<g<1.
\end{equation}
With $\tau$ in the tens or hundreds of milliseconds you hear a discrete
\textbf{echo} (a slap-back). With $\tau$ of only a few milliseconds your ear
cannot separate the two copies in time --- instead you hear a change in
\emph{timbre}. Why? Look at the frequency response. The impulse response is
$h[n]=\delta[n]+g\,\delta[n-D]$, so
\begin{equation}
\label{ch28:comb-H}
H(e^{\jj\omega})=1+g\,e^{-\jj\omega D},
\qquad
\abs{H(e^{\jj\omega})}^2 = 1+g^2+2g\cos(\omega D).
\end{equation}
The magnitude rises and falls cosinusoidally with frequency: peaks where the
delayed copy arrives in phase ($\omega D$ a multiple of $2\pi$) and dips where it
arrives out of phase. Plotted against frequency it looks like the teeth of a
comb --- hence \textbf{comb filter}. The teeth sit at multiples of $f_s/D=1/\tau$
hertz.

\begin{figure}[H]
\centering
\begin{tikzpicture}
\begin{axis}[
  width=12.5cm, height=5cm,
  xlabel={frequency $f$ (Hz), with $\tau=2$ ms so teeth every 500 Hz},
  ylabel={$|H(f)|$},
  xmin=0, xmax=4000, ymin=0, ymax=2.1,
  xtick={0,500,1000,1500,2000,2500,3000,3500,4000},
  xticklabel style={font=\scriptsize},
  ytick={0,0.5,1,1.5,2},
  axis lines=left, ymajorgrids, grid style={gray!18},
  samples=200, title={\small feedforward comb, $g=0.9$},
]
\addplot[domain=0:4000, thick, sigblue] {sqrt(1+0.81+2*0.9*cos(deg(2*pi*x*0.002)))};
\end{axis}
\end{tikzpicture}
\caption{The magnitude response of a single-echo (feedforward comb) filter,
Eq.~\eqref{ch28:comb-H}. Evenly spaced peaks and notches --- the comb teeth ---
spaced $1/\tau$ apart. A short delay is therefore a timbral effect, not a rhythmic
one: it colors the spectrum rather than producing an audible repeat.}
\label{ch28:fig-comb}
\end{figure}

\subsection{Repeating echoes and the feedback comb}

A real echo in a canyon repeats and fades. To get that, feed the output back
into the delay:
\begin{equation}
\label{ch28:feedback-comb}
y[n]=x[n]+g\,y[n-D].
\end{equation}
Now each output sample re-enters the delay line, so a single input impulse comes
back at $D, 2D, 3D,\dots$ with amplitudes $g, g^2, g^3,\dots$ --- a geometric
train of decaying repeats. This is an IIR filter; its impulse response is
infinite (it never exactly reaches zero), which is exactly what we want for a
decaying tail. Stability demands $\abs{g}<1$, or the echoes would grow without
bound and the system would scream.

\begin{walkbox}
How long does the feedback echo take to fade? Each pass multiplies by $g$, so
after $m$ repeats the level is $g^m$. The standard ``RT60'' measure asks when the
sound has dropped $60$~dB, i.e.\ to $10^{-3}$ of its start. Solve $g^m=10^{-3}$:
$m=-3/\log_{10}g$ repeats, taking $T_{60}=m\,\tau$ seconds. For $g=0.7$ and
$\tau=0.1$~s: $m\approx 19$ repeats, $T_{60}\approx 1.9$~s. The decay time is set
jointly by the feedback gain and the delay length --- the two knobs on a delay
pedal.
\end{walkbox}

\begin{lstlisting}[language=Python,caption={Feedback echo with a delay buffer.}]
import numpy as np

def echo(x, fs, delay_s=0.25, g=0.6):
    D = int(delay_s * fs)
    y = np.zeros(len(x) + D)
    y[:len(x)] = x
    for n in range(D, len(y)):           # y[n] = x[n] + g*y[n-D]
        y[n] += g * y[n - D]
    return y
\end{lstlisting}

\subsection{The all-pass filter: delay without coloration}

The comb filters above change the magnitude spectrum --- they put audible teeth
in it. But for natural-sounding reverb we want to smear the sound in time
\emph{without} stamping a comb-shaped color on it. The tool for that is the
\textbf{all-pass filter}: it passes every frequency at full, equal magnitude
(hence the name) while scrambling the \emph{phase}, and thus the timing.
\begin{equation}
\label{ch28:allpass}
y[n]=-g\,x[n]+x[n-D]+g\,y[n-D],
\qquad
\abs{H(e^{\jj\omega})}=1\ \text{for all }\omega.
\end{equation}

\begin{whybox}
How can a filter delay things yet leave \emph{every} magnitude untouched? Its
transfer function has the special form
$H(z)=\dfrac{-g+z^{-D}}{1-g\,z^{-D}}$: the numerator coefficients are the
denominator coefficients reversed. Each pole at radius $r$ is mirrored by a zero
at radius $1/r$ along the same ray. On the unit circle the contributions of every
mirror pair cancel exactly in magnitude --- so $\abs{H}=1$ everywhere --- but
their \emph{phases} add, producing a frequency-dependent delay. The all-pass is
pure time-smear with zero spectral coloring, which is precisely the ingredient a
reverberator needs to build density without ringing.
\end{whybox}

\subsection{Reverb: stacking combs and all-passes}

\textbf{Reverberation} is what a room does: thousands of reflections off walls,
floor, and ceiling arrive in a dense, decaying wash. We cannot afford to compute
thousands of literal echoes, so the classic \textbf{Schroeder reverberator}
(1962) fakes the statistics of a room cheaply:
\begin{itemize}
  \item Several \emph{feedback comb} filters in parallel, each with a different
        (mutually prime) delay, build up a dense set of decaying echoes. Coprime
        delays keep their teeth from lining up into an audible pitch.
  \item A few \emph{all-pass} filters in series then smear that into a smooth
        tail without adding coloration, multiplying the echo density at each
        stage.
\end{itemize}
A pre-delay (a plain delay before the network) sets the perceived room size, and
a low-pass filter in each feedback loop makes high frequencies decay faster than
lows --- which is what real air and soft furnishings do, and what makes a reverb
sound like a room rather than a tin can.

\begin{figure}[H]
\centering
\begin{tikzpicture}[
  box/.style={draw=gray!60, rounded corners=2pt, minimum height=8mm, minimum width=15mm, font=\scriptsize, align=center},
  >=Stealth, node distance=6mm,
]
\node[box, fill=sigblue!12] (in) {input};
\node[box, fill=sigblue!8, right=12mm of in] (c1) {comb $D_1$};
\node[box, fill=sigblue!8, below=3mm of c1] (c2) {comb $D_2$};
\node[box, fill=sigblue!8, above=3mm of c1] (c0) {comb $D_0$};
\node[box, fill=sigblue!8, below=3mm of c2] (c3) {comb $D_3$};
\node[circle, draw=gray!60, inner sep=1pt, right=12mm of c1, font=\scriptsize] (sum) {$+$};
\node[box, fill=quotientorange!18, right=8mm of sum] (a1) {allpass};
\node[box, fill=quotientorange!18, right=5mm of a1] (a2) {allpass};
\node[box, fill=targetgreen!15, right=5mm of a2] (out) {output};
\foreach \c in {c0,c1,c2,c3}{ \draw[->] (in.east) -- (\c.west); \draw[->] (\c.east) -- (sum); }
\draw[->] (sum) -- (a1);
\draw[->] (a1) -- (a2);
\draw[->] (a2) -- (out);
\end{tikzpicture}
\caption{A Schroeder reverberator: parallel feedback combs (coprime delays) build
echo density, then series all-passes smear the result into a smooth, uncolored
tail. Add a low-pass in each comb's feedback path for realistic high-frequency
decay.}
\label{ch28:fig-reverb}
\end{figure}

\subsection{Flanger and chorus: modulating the delay length}

Take the short-delay feedforward comb of Eq.~\eqref{ch28:feedforward-comb} and
\emph{slowly sweep the delay} $D$ with a low-frequency oscillator (LFO). As $D$
changes, the comb teeth slide up and down the spectrum. That sweeping comb is a
\textbf{flanger} --- the jet-plane ``whoosh'' --- with delays around 1--10~ms and
usually some feedback to deepen the notches.

\textbf{Chorus} is the same idea with a longer base delay (15--35~ms), a deeper,
slower, often randomized LFO, and little or no feedback. The longer delay puts
the comb teeth so close together that you do not hear them as coloration;
instead, the constantly shifting delay slightly detunes the copy (Doppler-style
pitch wobble), and mixing it with the original mimics several players slightly out
of tune and out of time --- the lush ``ensemble'' sound.

\begin{table}[H]
\centering
\begin{tabular}{@{}llll@{}}
\toprule
Effect & Base delay & Feedback? & What you hear \\
\midrule
Slap echo   & 50--150 ms & light & one discrete repeat \\
Delay/echo  & 100--800 ms & yes & rhythmic repeats, fading \\
Reverb      & dense spread & combs + all-pass & a room/space \\
Flanger     & 1--10 ms, swept & yes & sweeping ``whoosh'' \\
Chorus      & 15--35 ms, swept & little & thick ``ensemble'' \\
\bottomrule
\end{tabular}
\caption{The delay-based effects are one circuit --- a delay line plus optional
feedback plus optional modulation --- tuned to different delay ranges.}
\label{ch28:tab-delay}
\end{table}

\begin{trapbox}
Sweeping a delay means reading the buffer at a \emph{fractional} sample position
$n-D(n)$, where $D(n)$ is not an integer. Rounding to the nearest sample produces
audible ``zipper'' clicks. Use linear (or higher-order) interpolation between the
two neighboring samples in the delay buffer. Fractional-delay interpolation is what
keeps modulated delay lines smooth and free of zipper noise.
\end{trapbox}

\begin{recoverybox}
\textbf{Pause and think}: Why does a delay of $2$\,ms sound like a hollow tone coloration, while a delay of $200$\,ms sounds like a distinct echo? The human auditory system cannot resolve acoustic events separated by less than about $30$--$50$\,ms in time (the perceptual integration window, or Haas effect). Below this threshold, multiple arrivals fuse into a single perceived sound with an altered frequency spectrum (the comb filter teeth). Above it, the brain perceives two distinct auditory events separated in time.
\end{recoverybox}

\section{Synthesis I: additive --- building timbre by stacking sinusoids}
\label{ch28:additive}

Now we turn from shaping existing sound to creating it. The most direct route
follows straight from Fourier: \emph{any} periodic tone is a sum of harmonics, so
to build a tone, add up sinusoids.
\begin{equation}
\label{ch28:additive-eq}
y[n]=\sum_{k=1}^{K} A_k(n)\,\sin\!\bigl(2\pi k f_0 n/f_s + \phi_k\bigr).
\end{equation}
The fundamental $f_0$ sets the pitch; the relative amplitudes $A_k$ set the
\textbf{timbre} --- the ``color'' that distinguishes a flute from a violin at the
same pitch. Make the $A_k(n)$ change over time and you capture how a real
instrument's spectrum evolves (the bright attack of a plucked string mellowing
as upper harmonics die first).

\begin{intuitionbox}
Additive synthesis is the inverse Fourier series used as an instrument. You are
literally drawing the spectrum bar by bar and letting the synth play it back. It
can in principle reproduce \emph{any} harmonic sound --- it is the most general
periodic-tone synthesizer there is. Its drawback is cost and control: a rich tone
may need dozens of partials, each with its own time-varying amplitude curve, so
you are steering dozens of knobs at once. The other synthesis methods are largely
clever shortcuts that get a rich, evolving spectrum from far fewer controls.
\end{intuitionbox}

\begin{figure}[H]
\centering
\begin{tikzpicture}
\begin{axis}[
  width=12.5cm, height=4.6cm,
  xlabel={$n$ (samples)}, ylabel={amplitude},
  xmin=0, xmax=100, ymin=-1.7, ymax=1.7,
  xtick={0,25,50,75,100}, ytick={-1,0,1},
  axis lines=left, ymajorgrids, grid style={gray!18},
  legend style={at={(0.99,0.99)},anchor=north east,font=\scriptsize,draw=gray!40,legend columns=3},
  samples=200,
]
\addplot[domain=0:100, gray, dotted] {sin(deg(2*pi*x/50))};
\addlegendentry{$f_0$}
\addplot[domain=0:100, sigblue!60, dashed] {0.5*sin(deg(2*pi*3*x/50))};
\addlegendentry{$3f_0$}
\addplot[domain=0:100, thick, specred] {sin(deg(2*pi*x/50)) + 0.5*sin(deg(2*pi*3*x/50)) + 0.33*sin(deg(2*pi*5*x/50)) + 0.25*sin(deg(2*pi*7*x/50))};
\addlegendentry{sum (odd harmonics)}
\end{axis}
\end{tikzpicture}
\caption{Additive synthesis of a square-ish tone: the fundamental plus odd
harmonics with amplitudes $1,\tfrac13,\tfrac15,\tfrac17,\dots$ As more partials
are added the sum approaches a square wave. Choosing the harmonic amplitudes is
choosing the timbre.}
\label{ch28:fig-additive}
\end{figure}

\section{Synthesis II: subtractive --- carving timbre out of a rich source}
\label{ch28:subtractive}

The dominant approach in classic analog synthesizers reverses the additive
strategy. Instead of building a spectrum up from nothing, start with a
harmonically \emph{rich} waveform and \emph{carve away} what you do not want with
a filter. The signal chain is the iconic synth voice:
\[
\boxed{\text{oscillator}} \rightarrow \boxed{\text{filter}} \rightarrow \boxed{\text{amplifier}}
\]
\begin{itemize}
  \item \textbf{Oscillator (VCO).} A waveform rich in harmonics: a sawtooth (all
        harmonics, $A_k\propto 1/k$), a square/pulse (odd harmonics; pulse width
        changes the recipe), or a triangle (odd harmonics, fast roll-off). The
        sawtooth, with its full harmonic comb, is the standard starting point.
  \item \textbf{Filter (VCF).} A resonant low-pass whose \emph{cutoff} decides
        how much of the harmonic richness survives, and whose \emph{resonance} (a
        pole pushed near the unit circle) adds a vocal-like peak at the cutoff.
        Sweep the cutoff and you get the classic ``wah''/filter-sweep sound.
  \item \textbf{Amplifier (VCA).} Applies the loudness envelope (next section).
\end{itemize}

\begin{happeningbox}
The reason subtractive synthesis feels so expressive is that a single knob --- the
filter cutoff --- moves a huge amount of spectral energy at once. Opening the
filter as a note begins, then closing it as the note sustains, mimics exactly how
real plucked and struck instruments behave: bright at the attack, mellow as it
rings. One time-varying cutoff stands in for dozens of individual harmonic
envelopes that additive synthesis would force you to draw by hand. Subtractive
synthesis is additive synthesis with the harmonic envelopes \emph{ganged
together} under one control.
\end{happeningbox}

\begin{trapbox}
A naive sawtooth or square generated by $((n f_0/f_s)\bmod 1)$ has perfect, sharp
corners --- and therefore harmonics extending infinitely high, far past Nyquist.
Sampled directly, those fold back as aliasing and the oscillator sounds gritty
and detuned on high notes. Real synths use \emph{band-limited} oscillators
(BLEP/BLIT methods, or band-limited wavetables) that round the corners just
enough to keep every harmonic below $f_s/2$.
\end{trapbox}

\begin{figure}[H]
\centering
\begin{tikzpicture}
\begin{axis}[
  width=12.5cm, height=4.6cm,
  xlabel={frequency (harmonic number $k$)}, ylabel={amplitude},
  xmin=0, xmax=11, ymin=0, ymax=1.15,
  xtick={1,2,3,4,5,6,7,8,9,10}, ytick={0,0.5,1},
  axis lines=left, ymajorgrids, grid style={gray!18},
  legend style={at={(0.99,0.99)},anchor=north east,font=\scriptsize,draw=gray!40},
  title={\small a sawtooth's harmonics, before and after a low-pass at $k\!=\!4$},
]
\addplot+[ycomb, mark=*, thick, color=sigblue, mark options={fill=sigblue}]
  coordinates {(1,1)(2,0.5)(3,0.33)(4,0.25)(5,0.2)(6,0.167)(7,0.143)(8,0.125)(9,0.111)(10,0.1)};
\addlegendentry{sawtooth}
\addplot+[ycomb, mark=square*, thick, color=specred, mark options={fill=specred}]
  coordinates {(1,1)(2,0.5)(3,0.33)(4,0.20)(5,0.06)(6,0.02)(7,0.006)(8,0.002)(9,0.0008)(10,0.0003)};
\addlegendentry{after low-pass}
\end{axis}
\end{tikzpicture}
\caption{Subtractive synthesis in the spectrum. Blue: the sawtooth's full harmonic
series, $A_k=1/k$. Red: after a resonant low-pass with cutoff near the 4th
harmonic --- harmonics above the cutoff are carved away (with a small resonant
bump right at the cutoff). Sweeping the cutoff slides the red ``cliff'' up and
down.}
\label{ch28:fig-subtractive}
\end{figure}

\section{Synthesis III: FM --- rich timbre from a wobbling phase}
\label{ch28:fm}

Here is where nonlinearity returns in spectacular fashion. \textbf{Frequency
modulation} synthesis, discovered by John Chowning at Stanford in the late 1960s,
gets startlingly complex, evolving spectra from just \emph{two} oscillators by
using one to wobble the other's frequency:
\begin{equation}
\label{ch28:fm-eq}
y(t)=\sin\!\Bigl(2\pi f_c t + I\sin(2\pi f_m t)\Bigr).
\end{equation}
The \emph{carrier} at $f_c$ has its instantaneous phase shoved back and forth by
the \emph{modulator} at $f_m$. The single most important number is the
\textbf{modulation index} $I$: it controls how violently the phase wobbles, and
therefore how rich the spectrum becomes.

\begin{historybox}
John Chowning stumbled onto digital FM synthesis in 1967 at Stanford's CCRMA while
experimenting with extreme vibrato: as he pushed the vibrato modulation rate beyond $20$\,Hz
into the audible range, the pitch flutter fused into a completely new, harmonically rich
\emph{timbre}. At a time when mainframe computers took hours of computation to sum the
dozens of sine waves required for additive synthesis, Chowning realized that just two
interacting oscillators could produce a rich spectrum of dynamic overtones via Bessel
function sideband distribution. Stanford patented the discovery in 1973, but American
analog synthesizer companies initially dismissed digital FM as an unnatural, inharmonic
mathematical curiosity. Yamaha licensed the patent exclusively and released the DX7 in 1983.
The DX7 sold over $200{,}000$ units,
democratized digital synthesis, and defined the sonic signature of 1980s music, while
compact Yamaha OPL chips powered audio across a generation of personal computers. An active
modern research challenge is inverse FM synthesis: using differentiable DSP and neural
optimization to automatically discover the operator topologies, frequency ratios, and
modulation envelopes that reproduce an arbitrary real-world acoustic instrument recording.
\end{historybox}

\begin{whybox}
Why does wobbling the phase create a whole comb of sidebands? Expanding
Eq.~\eqref{ch28:fm-eq} with the Jacobi--Anger identity gives
\[
y(t)=\sum_{k=-\infty}^{\infty} J_k(I)\,\sin\!\bigl(2\pi(f_c+k f_m)t\bigr),
\]
where $J_k$ are Bessel functions. The spectrum is a \emph{ladder} of components
at $f_c, f_c\pm f_m, f_c\pm 2f_m, \dots$, with strengths $J_k(I)$
(Figure~\ref{ch28:fig-fm-spectrum}). When $I$ is
small only $k=0,\pm1$ matter and you hear a nearly pure tone. As $I$ grows, the
energy spreads out to higher $|k|$ --- the timbre brightens and fills in.
\emph{One knob, $I$, sweeps from a sine to a brassy blare.} That is why a single
FM ``operator pair'' driven by an envelope on $I$ can imitate brass, bells, and
electric pianos so convincingly.
\end{whybox}

\begin{aestheticbox}
The ratio $f_m/f_c$ decides whether FM sounds musical or clangorous. When the
ratio is a simple integer (1:1, 2:1, 3:2), the sidebands $f_c+k f_m$ land on
exact harmonics of a common fundamental, and you hear a clear pitch with a rich
but harmonic timbre. When the ratio is irrational or complex, the sidebands fall
between harmonics, giving the inharmonic, metallic, bell-like tones FM is famous
for. Two oscillators and two numbers --- $I$ and the ratio --- span an enormous
expressive range. This efficiency is exactly why FM ruled 1980s digital
synthesizers, where memory and compute were precious.
\end{aestheticbox}

\begin{figure}[H]
\centering
\begin{tikzpicture}
\begin{axis}[
  width=12.5cm, height=4.6cm,
  xlabel={$t$}, ylabel={amplitude},
  xmin=0, xmax=1, ymin=-1.4, ymax=1.4,
  xtick={0,0.25,0.5,0.75,1}, ytick={-1,0,1},
  axis lines=left, ymajorgrids, grid style={gray!18},
  samples=200, title={\small FM: carrier $f_c\!=\!8$, modulator $f_m\!=\!2$, index $I\!=\!4$},
]
\addplot[domain=0:1, thick, sigblue] {sin(deg(2*pi*8*x + 4*sin(deg(2*pi*2*x))))};
\end{axis}
\end{tikzpicture}
\caption{An FM waveform. Notice how the carrier's local frequency speeds up and
slows down --- the cycles bunch and stretch --- as the modulator pushes its
phase. In the spectrum this becomes a ladder of sidebands whose richness is set
by the modulation index $I$.}
\label{ch28:fig-fm}
\end{figure}

\begin{figure}[H]
\centering
\begin{tikzpicture}
\begin{axis}[
  width=12.5cm, height=4.6cm,
  xlabel={sideband index $k$ (component at $f_c+k f_m$)},
  ylabel={$|J_k(I)|$},
  xmin=-8, xmax=8, ymin=0, ymax=0.52,
  xtick={-7,-6,-5,-4,-3,-2,-1,0,1,2,3,4,5,6,7},
  xticklabel style={font=\scriptsize},
  ytick={0,0.25,0.5},
  axis lines=left, ymajorgrids, grid style={gray!18},
  title={\small FM sideband strengths at modulation index $I=4$},
]
\addplot+[ycomb, mark=*, thick, color=sigblue, mark options={fill=sigblue}]
  coordinates {(-7,0.015)(-6,0.049)(-5,0.132)(-4,0.281)(-3,0.430)(-2,0.364)(-1,0.066)(0,0.397)(1,0.066)(2,0.364)(3,0.430)(4,0.281)(5,0.132)(6,0.049)(7,0.015)};
\end{axis}
\end{tikzpicture}
\caption{The FM spectrum for index $I=4$: component $k$ sits at $f_c+k f_m$ with
amplitude $|J_k(I)|$ given by the Jacobi--Anger Bessel expansion above. Notice the carrier
($k=0$) is no longer dominant --- energy has spread out into the sidebands. As
$I$ grows the ladder reaches farther out and the tone brightens; as $I\to 0$ only
$k=0,\pm 1$ survive and the tone returns to a near-pure sine.}
\label{ch28:fig-fm-spectrum}
\end{figure}

\begin{lstlisting}[language=Python,caption={Two-operator FM with an envelope on the index.}]
import numpy as np

def fm_tone(fc, fm, dur, fs=44100, I0=5.0, tau=0.4):
    t = np.arange(int(dur*fs)) / fs
    I = I0 * np.exp(-t / tau)            # index decays -> bright attack, mellow tail
    return np.sin(2*np.pi*fc*t + I*np.sin(2*np.pi*fm*t))

bell = fm_tone(fc=440, fm=440*1.41, dur=2.0)   # near-irrational ratio -> metallic
\end{lstlisting}

\section{Synthesis IV: wavetable --- play back one stored cycle}
\label{ch28:wavetable}

The cheapest way to get an arbitrary timbre is to \emph{store} one period of the
desired waveform in a table and read it out in a loop. A \textbf{phase
accumulator} steps an index through the table; how fast it steps sets the pitch.
\begin{equation}
\label{ch28:phase-acc}
\phi[n]=\bigl(\phi[n-1]+\Delta\bigr)\bmod N,
\qquad
\Delta = \frac{f_0\,N}{f_s},
\qquad
y[n]=\text{table}\bigl[\phi[n]\bigr].
\end{equation}
Here $N$ is the table length and $\Delta$ the phase increment. A larger increment
skips through the table faster, raising the pitch. Because $\Delta$ is generally
fractional, you again interpolate between table entries.

\begin{intuitionbox}
A wavetable oscillator is just additive synthesis \emph{precomputed} and frozen:
sum your harmonics once, store the resulting cycle, then replay it cheaply
forever. You trade the live control of additive synthesis for almost no runtime
cost. ``Wavetable synthesis'' in modern soft-synths (Serum, Massive) extends this
by storing \emph{many} tables and smoothly crossfading between them as a note
evolves --- morphing the timbre in time, which is exactly the evolving spectrum
that additive synthesis offered, recovered through a stack of snapshots.
\end{intuitionbox}

\begin{trapbox}
A single stored cycle has a fixed harmonic content, so transposing it far upward
will push its upper harmonics past Nyquist and alias. Production wavetable
synths keep a set of band-limited versions of each table (\emph{mip-maps}), one
per octave, and pick the version whose highest harmonic still fits below $f_s/2$
for the note being played.
\end{trapbox}

\section{Synthesis V: Karplus--Strong --- physical modeling from a noisy loop}
\label{ch28:karplus-strong}

Additive, subtractive, FM, and wavetable synthesis are abstract mathematical recipes.
\textbf{Physical modeling synthesis}, by contrast, builds a discrete-time simulator of an
actual vibrating object. The classic example in this family is the \textbf{Karplus--Strong algorithm}
(Kevin Karplus and Alexander Strong, 1983), which produces a remarkably realistic acoustic guitar
or harp pluck using nothing more than a delay buffer, an adder, and a two-tap averager.

\subsection{The mechanism: pluck with noise, sustain with feedback}

Picture a guitar string. When a plectrum plucks it, it delivers a brief burst of chaotic,
wideband kinetic energy (white noise). That displacement wave bounces back and forth between
the fixed bridge and nut. On every round trip, internal friction and air drag slightly blur
the sharp wave peaks and dissipate high frequencies faster than low ones.

Karplus--Strong simulates this directly:
\begin{enumerate}
  \item \textbf{Excitation}: Fill a delay line of length $N$ with $N$ samples of random white noise.
        This represents the initial displacement profile along the plucked string.
  \item \textbf{Recirculation}: Read the oldest sample from the delay line, pass it through an ultra-simple
        one-pole low-pass filter (averaging adjacent samples), attenuate by a decay factor $\alpha \lesssim 1$,
        and write the result back into the front of the delay line.
  \item \textbf{Output}: Stream the circulating samples directly to the DAC.
\end{enumerate}

\begin{equation}
\label{ch28:ks-eq}
y[n] = x[n] + \alpha \cdot \frac{y[n-N] + y[n-(N+1)]}{2},
\end{equation}
where $x[n]$ is the initial noise burst (non-zero only for $0 \le n < N$).

The round-trip delay $N$ sets the fundamental pitch of the note:
\begin{equation}
\label{ch28:ks-pitch}
f_0 \approx \frac{f_s}{N + \tfrac{1}{2}},
\end{equation}
where the extra $\tfrac{1}{2}$ sample accounts for the phase delay introduced by the two-point averager.
As the waveform circulates hundreds of times per second, the low-pass filter progressively shaves off
the high harmonics. The harsh noise burst dissolves within milliseconds into a pure, ringing, warm
harmonic tone that naturally decays into silence.

\begin{figure}[H]
\centering
\begin{tikzpicture}[font=\small, >=Stealth]
\tikzset{
  blk/.style={draw=blue!70!black, fill=sourceblue!10, rounded corners=2pt, minimum height=9mm, minimum width=22mm, align=center, thick},
  flt/.style={draw=targetgreen!70!black, fill=targetgreen!10, rounded corners=2pt, minimum height=9mm, minimum width=24mm, align=center, thick},
  gain/.style={draw=quotientorange!80!black, fill=quotientorange!12, rounded corners=2pt, minimum height=9mm, minimum width=16mm, align=center, thick},
  sum/.style={draw=gray!70!black, fill=gray!10, circle, inner sep=2pt, thick}
}

% Nodes
\node (in) at (0,0) [font=\normalsize] {Burst $x[n]$};
\node[sum] (sum) at (2.4,0) {\small $+$};
\node[blk] (delay) at (5.4,0) {Delay Line\\$z^{-N}$ ($N$ samples)};
\node (out) at (9.4,0) [font=\normalsize, sigblue] {$y[n]$ (Pluck)};

\node[flt] (lpf) at (6.6, -1.8) {Lowpass Averager\\$\frac{1 + z^{-1}}{2}$};
\node[gain] (decay) at (3.8, -1.8) {Decay $\alpha$\\($\alpha \approx 0.99$)};

% Forward path
\draw[->, thick] (in) -- (sum);
\draw[->, thick] (sum) -- (delay);
\draw[->, thick] (delay) -- (out);

% Feedback path
\draw[->, thick] (7.8,0) |- (lpf.east);
\draw[->, thick] (lpf) -- (decay);
\draw[->, thick] (decay) -| (sum.south) node[pos=0.8, left, font=\scriptsize] {$+$};

\end{tikzpicture}
\caption{The Karplus--Strong plucked-string synthesis loop. An initial noise burst of $N$ samples circulates through an $N$-sample delay line. On each round trip, a two-point averager dampens high frequencies while a loss factor $\alpha$ dissipates overall energy. The loop circulation period $N$ dictates the fundamental frequency $f_0 \approx f_s / (N + 0.5)$, producing a convincing guitar pluck with physically accurate harmonic decay.}
\label{ch28:fig-karplus-strong}
\end{figure}

\begin{lstlisting}[language=Python,caption={The Karplus--Strong string synthesis algorithm in Python.}]
import numpy as np

def karplus_strong(f0, dur, fs=44100, alpha=0.996):
    N = int(round(fs / f0 - 0.5))
    buf = np.random.uniform(-1.0, 1.0, N)       # initial white noise excitation
    total_samples = int(dur * fs)
    y = np.zeros(total_samples)
    y[:N] = buf
    
    # Recirculation loop with lowpass averaging
    for i in range(N, total_samples):
        val = alpha * 0.5 * (buf[0] + buf[1])
        y[i] = val
        buf = np.roll(buf, -1)                  # advance the delay line
        buf[-1] = val
    return y
\end{lstlisting}

\section{The envelope: shaping a sound in time (ADSR)}
\label{ch28:adsr}

An oscillator alone drones forever at constant loudness --- nothing like a real
note, which is born, swells, settles, and dies. The \textbf{amplitude envelope}
is the time-domain shape multiplied onto the sound, and the universal model for
it is \textbf{ADSR}:
\begin{itemize}
  \item \textbf{Attack} --- time to rise from silence to peak after the note
        starts (a key press, or a wrist flick).
  \item \textbf{Decay} --- time to fall from the peak to the sustain level.
  \item \textbf{Sustain} --- the held \emph{level} (not a time) maintained while
        the note is held.
  \item \textbf{Release} --- time to fall back to silence after the note ends.
\end{itemize}

\begin{figure}[H]
\centering
\begin{tikzpicture}
\begin{axis}[
  width=12.5cm, height=4.6cm,
  xlabel={time}, ylabel={amplitude},
  xmin=0, xmax=10, ymin=0, ymax=1.15,
  xtick=\empty, ytick={0,0.6,1}, yticklabels={$0$,sustain,peak},
  axis lines=left, ymajorgrids, grid style={gray!18},
]
\addplot[thick, sigblue] coordinates {(0,0)(1.5,1)(3,0.6)(6.5,0.6)(8.5,0)};
\draw[gray!55, dashed] (axis cs:1.5,0)--(axis cs:1.5,1);
\draw[gray!55, dashed] (axis cs:3,0)--(axis cs:3,0.6);
\draw[gray!55, dashed] (axis cs:6.5,0)--(axis cs:6.5,0.6);
\node[font=\scriptsize] at (axis cs:0.75,1.07) {A};
\node[font=\scriptsize] at (axis cs:2.25,1.07) {D};
\node[font=\scriptsize] at (axis cs:4.75,1.07) {S (held)};
\node[font=\scriptsize] at (axis cs:7.5,1.07) {R};
\node[font=\scriptsize, gray] at (axis cs:6.5,0.08) {note off};
\end{axis}
\end{tikzpicture}
\caption{The ADSR envelope. Attack and decay fire when the note starts; the
sustain level holds while the key is down; release begins when the note is let
go. Multiplying the oscillator by this curve turns a drone into a note.}
\label{ch28:fig-adsr}
\end{figure}

\begin{happeningbox}
The envelope is most of what your ear uses to name an instrument. Classic
experiments show that if you splice off the attack transient of a recorded piano,
trumpet, and violin and play only their steady sustains, listeners struggle to
tell them apart --- the \emph{attack} carries the identity. A fast, percussive
attack with no sustain says ``plucked'' or ``struck''; a slow attack with a long
sustain says ``bowed'' or ``blown.'' This is why the same oscillator can become a
plucked bass or a swelling pad purely by changing the envelope. Timbre lives in
the spectrum \emph{and} its evolution; the envelope governs the evolution.
\end{happeningbox}

Envelopes do not have to drive loudness only. Route an ADSR (or a second one) to
the filter cutoff and you get the bright-attack-then-mellow behavior of
subtractive synthesis; route it to the FM modulation index and you get the brassy
swell of an FM trumpet. The envelope is a general-purpose, one-shot control
signal --- and that makes it the natural place for a gesture to take over.

\begin{lstlisting}[language=Python,caption={A linear ADSR envelope generator.}]
import numpy as np

def adsr(fs, a, d, s_level, s_time, r):
    seg = lambda t, lo, hi: np.linspace(lo, hi, int(t*fs), endpoint=False)
    return np.concatenate([
        seg(a, 0.0, 1.0),               # attack
        seg(d, 1.0, s_level),           # decay
        seg(s_time, s_level, s_level),  # sustain (held)
        seg(r, s_level, 0.0),           # release
    ])

# voice = oscillator(...) * adsr(44100, 0.01, 0.1, 0.6, 0.5, 0.3)
\end{lstlisting}

\section{Putting it together: a complete synth voice}
\label{ch28:voice}

We can now read the block diagram of essentially any monophonic synthesizer
voice as a sentence in the language of this book.

\begin{figure}[H]
\centering
\begin{tikzpicture}[
  blk/.style={draw=gray!60, rounded corners=2pt, minimum height=9mm, minimum width=16mm, font=\scriptsize, align=center},
  ctl/.style={draw=quotientorange!70, rounded corners=2pt, minimum height=8mm, minimum width=18mm, font=\scriptsize, align=center, fill=quotientorange!10},
  >=Stealth,
]
\node[blk, fill=sigblue!10] (osc) {OSC\\(rich wave)};
\node[blk, fill=sigblue!10, right=12mm of osc] (vcf) {VCF\\(low-pass)};
\node[blk, fill=sigblue!10, right=12mm of vcf] (vca) {VCA\\(gain)};
\node[blk, fill=targetgreen!15, right=10mm of vca] (out) {output};
\draw[->] (osc) -- (vcf);
\draw[->] (vcf) -- (vca);
\draw[->] (vca) -- (out);
\node[ctl, below=10mm of vcf] (env1) {ADSR $\to$ cutoff};
\node[ctl, below=10mm of vca] (env2) {ADSR $\to$ gain};
\node[ctl, below=10mm of osc] (lfo) {LFO / pitch};
\draw[->, quotientorange!80] (env1) -- (vcf);
\draw[->, quotientorange!80] (env2) -- (vca);
\draw[->, quotientorange!80] (lfo) -- (osc);
\end{tikzpicture}
\caption{A canonical subtractive voice. Blue: the audio path (oscillator
$\to$ filter $\to$ amplifier). Orange: the control signals --- envelopes and an
LFO --- that modulate the audio path over time. Every electronic instrument is a
variation on this skeleton; the gesture controller's job is to supply or replace
the orange control signals.}
\label{ch28:fig-voice}
\end{figure}

\begin{applicationbox}
\textbf{Unifying the synthesis palette.} A drum machine is this voice with a noise
oscillator and a very short envelope. A pad is this voice with a slow attack and
long release. An organ is several oscillators (additive) with a near-square
envelope. A DX-style electric piano is the FM voice of
Section~\ref{ch28:fm} with a decaying index envelope. Stack a delay, a chorus,
and a reverb after the VCA and you have a finished, recordable sound. There is
nothing in an expensive synthesizer that is not a tuned, polished arrangement of
the boxes in this chapter.
\end{applicationbox}

\section{Toward the wrist tracker: gesture as a control signal}
\label{ch28:goal}

\begin{goalbox}
Everything in this chapter splits cleanly into two streams: the \emph{audio path}
(oscillators, filters, delays --- the blue boxes) and the \emph{control path}
(envelopes, LFOs, knobs --- the orange boxes). The audio path runs at the sample
rate, $44.1$~kHz, far faster than any human motion. The control path runs slowly,
tens to a few hundred updates per second --- and that is exactly the rate at which
a wrist sensor produces useful information.

So the design is clean: \textbf{let the wrist drive the control path}, and let a
software synth (or a hardware one over MIDI) own the audio path. Concretely:
\begin{itemize}
  \item A coarse, discrete \emph{gesture} (a flick, a fist, a wrist-tilt held)
        is a reliable trigger --- map it to MIDI \emph{note-on/note-off}, the
        moment the ADSR fires.
  \item A continuous, smoothly varying \emph{orientation} (wrist pitch or roll
        from the fused IMU of Chapter~36) is a reliable knob --- map
        it to a MIDI \emph{continuous controller}: filter cutoff, FM index,
        vibrato depth, or volume. Slow-moving, low-dimensional, exactly what the
        control path wants.
\end{itemize}
This is why the synthesis architecture matters to the build: it tells us we do
\emph{not} need to recover anything fast or finely detailed to make expressive
music. A handful of clean, slow control streams modulating a well-designed voice
is already a playable instrument.
\end{goalbox}

\begin{trapbox}
\textbf{Know what the wrist can supply.} A wrist-worn IMU measures the
\emph{orientation and motion of the whole forearm/wrist}, not the angle of each
finger. It cannot recover continuous individual finger-joint angles --- the
fingers are mechanically downstream of the wrist and largely invisible to it.
What it gives reliably is (a) coarse discrete gestures (flick, twist, tap,
clench-versus-relax via the accompanying jolt) and (b) smooth wrist/forearm
orientation. Individual-finger control needs sensors that see the muscles or the
fingers themselves --- surface EMG or force myography on the forearm for finger
\emph{activation}, or instrumented rings for finger \emph{pose} (Chapter~37
goes into this). Designing the instrument around triggers and slow knobs
is not a compromise forced by laziness; it is matching the controller to what the
physics of the sensor can actually deliver.
\end{trapbox}

\begin{buildbox}
A minimal end-to-end test, before any custom hardware: run a small synth on a
laptop or Raspberry Pi (e.g.\ a SuperCollider, Pure Data, or a 50-line NumPy
voice), and feed it MIDI from a microcontroller on the wrist.

\textbf{Firmware side (Arduino-style pseudocode):}
\begin{lstlisting}[language=C,caption={Wrist control to MIDI: trigger + one continuous knob.}]
// Runs at ~100 Hz. imu gives fused pitch (radians) and accel-based jerk.
float pitch = imu_pitch();           // smooth wrist tilt: a "knob"
float jerk  = imu_accel_jerk();      // sudden motion: a "trigger"

// Continuous controller: map wrist tilt [-pi/4, pi/4] -> CC 74 (filter cutoff)
int cc = clamp((int)((pitch + 0.785f) / 1.571f * 127.0f), 0, 127);
midiControlChange(0, 74, cc);        // channel 0, CC #74

// Discrete trigger: a sharp flick fires a note (with simple debounce)
if (jerk > THRESH && (now - last_note) > 120 /*ms*/) {
    midiNoteOn(0, 60, 100);          // middle C, velocity 100
    last_note = now;  note_active = true;
} else if (note_active && jerk < REL_THRESH) {
    midiNoteOff(0, 60, 0);  note_active = false;
}
\end{lstlisting}

Notes from experience: smooth the IMU-derived knob with a one-pole low-pass or the
MIDI value will jitter and the filter will buzz; debounce the trigger (a single
flick produces a burst of high-acceleration samples); and send CC updates only
when the value actually changes, to avoid flooding the MIDI bus.
\end{buildbox}

\section{Summary}
\label{ch28:summary}

Audio effects and synthesis are not a new theory --- they are the filters,
delays, and oscillators of the earlier chapters, arranged with intent.

\begin{itemize}
  \item Every effect turns one of three knobs: \textbf{amplitude} (gain,
        tremolo, distortion), \textbf{spectrum} (EQ = filtering), or
        \textbf{time} (delay). Movement comes from \emph{modulating} a knob.
  \item \textbf{Distortion} is a memoryless nonlinearity; bending the amplitude
        curve sprays the spectrum into harmonics. Watch out for aliasing.
  \item \textbf{Delay} seeds the spatial effects. A short delay is a \emph{comb
        filter} (timbral); feedback makes it a decaying echo; \emph{all-pass}
        filters smear time without coloring the spectrum; stacking combs and
        all-passes makes \emph{reverb}; modulating a short delay makes
        \emph{flanger} and \emph{chorus}.
  \item \textbf{Synthesis} builds timbre five ways: \emph{additive} (sum
        sinusoids), \emph{subtractive} (filter a rich source), \emph{FM} (modulate
        phase to generate Bessel sidebands from two oscillators), \emph{wavetable}
        (replay a stored cycle), and \emph{physical modeling} (simulate a resonating
        waveguide such as Karplus--Strong).
  \item The \textbf{ADSR envelope} shapes a sound in time and carries much of an
        instrument's identity, especially in its attack.
  \item For the wrist controller, the split between a fast \emph{audio path} and
        a slow \emph{control path} is the design key: drive the control path with
        discrete gesture triggers (note-on/off) and smooth orientation knobs
        (continuous controllers), matched to what a wrist IMU can actually sense.
\end{itemize}

Having built the core blocks of digital effects and synthesis, we turn in the next
chapter from synthetic algorithms to the acoustic physics of real musical instruments.
In Chapter~29, we examine how vibrating strings, air columns, percussion modes, and
the MIDI protocol dictate the latency and harmonic structure of expressive digital instruments.

\section{Exercises}
\label{ch28:exercises}

\begin{exercise}
A pure tone $x[n]=\cos(\omega_0 n)$ is fed through the cubic waveshaper
$f(x)=x-\tfrac13 x^3$. Using $\cos^3\theta=\tfrac34\cos\theta+\tfrac14\cos3\theta$,
find the frequencies present in the output and their amplitudes. Which new
frequency appears, and is it a harmonic of $\omega_0$?
\end{exercise}

\begin{exercise}
A feedforward comb filter has $h[n]=\delta[n]+g\,\delta[n-D]$ with $f_s=48$~kHz
and $D=24$ samples. (a) What is the delay $\tau$ in milliseconds? (b) At what
frequencies (in Hz) are the comb's first three peaks? (c) If $g=0.95$, what are
the maximum and minimum values of $\abs{H(e^{\jj\omega})}$?
\end{exercise}

\begin{exercise}
For the feedback echo $y[n]=x[n]+g\,y[n-D]$ with $\tau=300$~ms, choose $g$ so the
echoes decay by $60$~dB in about $3$~seconds. (Use $g^m=10^{-3}$ with
$m=T_{60}/\tau$.) Then state the danger if $g$ is set to exactly $1.0$.
\end{exercise}

\begin{exercise}
Verify the all-pass magnitude claim for $D=1$:
$H(z)=\dfrac{-g+z^{-1}}{1-g z^{-1}}$. Substitute $z=e^{\jj\omega}$ and show
algebraically that $\abs{H(e^{\jj\omega})}=1$ for all $\omega$ (numerator and
denominator have equal magnitude). What does this tell you about why all-pass
filters are used to build reverb tails?
\end{exercise}

\begin{exercise}
You want to additively synthesize a sawtooth at $f_0=200$~Hz, sampled at
$f_s=44.1$~kHz, using harmonics $A_k=1/k$. What is the largest harmonic index
$k$ you may include before exceeding Nyquist? What goes wrong audibly if you
include harmonics beyond that?
\end{exercise}

\begin{exercise}
An FM voice has $f_c=300$~Hz and modulator $f_m=300$~Hz (ratio $1{:}1$). (a) List
the frequencies of the spectral components for $k=-2,\dots,2$. (b) Are they all
harmonics of a common fundamental? (c) Now set $f_m=300\sqrt{2}$~Hz; describe
qualitatively how the timbre changes and why.
\end{exercise}

\begin{exercise}
A wavetable of length $N=2048$ is played at $f_s=44.1$~kHz. (a) What phase
increment $\Delta$ produces $f_0=440$~Hz? (b) Since $\Delta$ is not an integer,
what must you do when reading the table, and what artifact appears if you simply
round to the nearest index?
\end{exercise}

\begin{exercise}
Sketch (or compute in code) two ADSR envelopes: one for a plucked string and one
for a bowed pad. Give rough numbers (in ms and sustain level) for each of A, D,
S, R, and justify each choice in terms of how the real instrument behaves.
\end{exercise}

\begin{exercise}
Design a mapping from a wrist IMU to a subtractive synth voice. Specify: which
sensor quantity triggers note-on/note-off, which quantity drives the filter
cutoff (a MIDI continuous controller), and how you would smooth and debounce
each. Then state explicitly one musical gesture this scheme \emph{cannot}
support, and which additional sensor (EMG, FMG, or a finger ring) would be needed
to support it.
\end{exercise}

\begin{exercise}
\textbf{(Code)} Implement a chorus effect: a delay line of base length 25~ms
whose read position is modulated by a $0.5$~Hz sine of depth $\pm 5$~ms, mixed
$50/50$ with the dry signal, using linear interpolation for the fractional delay.
Run it on a sustained synth note and on a drum loop. Describe how the effect
differs on the two sources and why (think about what a slowly varying delay does
to a steady tone versus to sharp transients).
\end{exercise}

\begin{exercise}
A tremolo is $y[n]=(1+d\cos(2\pi f_m n/f_s))\,x[n]$ applied to a tone
$x[n]=\cos(2\pi f_0 n/f_s)$. Use a product-to-sum identity to show the output
contains components at $f_0$ and $f_0\pm f_m$. Why, at $f_m\approx 5$~Hz, do you
perceive a single pulsing tone rather than three separate pitches?
\end{exercise}

\begin{exercise}
\textbf{(Concept)} Explain in two or three sentences why distortion, FM
synthesis, and the harmonics of a clipped oscillator are all manifestations of
the \emph{same} underlying phenomenon, and contrast that with EQ and echo, which
share a different underlying phenomenon. Name the two phenomena.
\end{exercise}
